\scanpage{15}
\mspage{\handwritten{29}}

\deleted{radicielle finie (par exemple, soit séparable)}. \inserted{\handwritten{Conditions \unsure{équivalentes}: 1)}} \deleted{Pour que} $\mathrm B$ \deleted{soit}\inserted{\handwritten{est}} une algèbre de Cohen sur $\mathrm A$ (sous-entendu: pour les topologies habituelles, - il faudrait le préciser dans \deleted{\unsure{3}}\handwritten{4}.7.) , \deleted{il f et suffit que}\inserted{\handwritten{2)}} $\mathrm B$ \deleted{\unsure{soit|est}}\handwritten{est} plat sur $\mathrm A$, complet, et \deleted{que} $\mathrm B_{\unsure{0|o}}=\mathrm B/\underline{\mathrm m}\mathrm B$ \deleted{soit}\inserted{\handwritten{est}} une algèbre de Cohen sur $k=\mathrm A/\underline{\mathrm m}$, i.e. formellement simple sur $k$, \inserted{\handwritten{cf 2.\unsure{6|5}.}} \deleted{\handwritten{3) $B$ est \inserted{\unsure{complet et}} formellement simple sur $A$ \unsure{à fibre|et tel que|ou plat et} \unsure{séparable}}}

(On aura noté \deleted{pa} après \deleted{3}\handwritten{4}.7. que pour une algèbre topologique sur un corps $k$, Cohen = adique et formellement simple ). \deleted{La nécessité}\inserted{\handwritten{1)$\Rightarrow$2) \deleted{\unsure{et 1)$\Rightarrow$3)}}}} étant déjà acquis\deleted{e} (sans hypothèse de multiplicité radicielle finie), prouvons \deleted{la suffisance}\inserted{\handwritten{2)$\Rightarrow$1)}}. En vertu de la définition, on est ramené simplement à prouver que $\mathrm B$ satisfait la condition $\mathrm P$. Cela nous ramène au cas où $\mathrm A$ est artinien. En vertu du N\textsuperscript{\uline{o}}2, on sait que $\mathrm B_0$ est régulier, considérons un système régulier de paramètres de $\mathrm B_0$ (un système de paramètres formant une $\mathrm B_0$-suite suffirait), relevons le en un système de paramètres de $\mathrm B$, soit $x=(x_1,\ldots,x_r)$. Alors les idéaux engendrés par les $x^n=(x_1^n,\ldotp\ldotp,x_r^n)$, pour $n\handwritten{\geq}0$, forment un système fondamental de voisinages \inserted{\handwritten{$\unsure{I_n|J_n}$}} de 0 \deleted{\handwritten{\unsure{de}}} dans $\mathrm B$, d'autre part on constate aussitôt que les $\mathrm B/\mathrm J_n$ sont des modules plats sur $\mathrm A$ (il y a un \deleted{lemme} corollaire dans l\deleted{\unsure{q}}a platitude supérieure, dans le cas où on divise par un seul élement \deleted{non divi} régulier, cas auquel on est ramené aussitôt). Comme $\mathrm A$ est artinien, ils sont donc projectifs. \deleted{\unsure{c.q.f.d.}} \deleted{\handwritten{\unsure{Enfin, on aura} 3)$\Rightarrow$2) \unsure{si $k(\mathrm B)/k$ séparable} \unsure{grâce à} 3.\unsure{8}. \unsure{(passé de 3) à 1))} en utilisant le \uline{th. d'existence} \uline{ci-dessous} 3.\unsure{9|11}.}}

\deleted{NB\quad On n'a utilisé l'hypothèse sur la multiplicité radicielle que pour \deleted{pouvoir} savoir que $\mathrm B_0$ est Cohen-Macaulay. L'hypothèse faite est très probablement inutile. Serre vient de me dire que Mac-Lane a prouvé la réciproque de 2.3. sans hypothèse de mu\deleted{\unsure{ll}}ltiplicité rad finie, dans ``Note on a relative structure of p-adic fi\deleted{\unsure{el}}elds'' Ann Math t.41 1940, page 751 -753. Cela vaudrait peut-être le coup, pour \deleted{liquider d} le coup d'oeuil, que tu regardes cela et donnes le résultat complet dans 2.3.}\margin{\handwritten{$\Vert$}} On
% A double vertical stroke stands beside the citation and page number. The ruled cancellation steps around the last On.
% archive sheet number: 14
